% =====================================================================
%  Relation to Torr's Bayesian framework, GRIC and MAGSAC++
%  (camera-ready; rebuttal to Reviewer P2b3, Question 2, and Reviewer
%  vKBs, "MAGSAC++ discussion"; the AC's Action 5).  \input at the end
%  of the "Shekhovtsov's Unified Threshold Framework" appendix.
% =====================================================================
\subsection{Relation to Torr's Bayesian Framework, \gric, and \magsacpp}
\label{sec:supp-torr}

\paragraph{Torr's Bayesian manifold formulation.}
Torr's Bayesian treatment of epipolar geometry and generic manifold
fitting~\cite{Torr2002} already shows that marginalizing the latent
noiseless correspondence removes the tangential error, leaving a
Gaussian likelihood in the manifold codimension~$d_g$, together with a
uniform outlier model.  This is the foundation our
likelihood~\eqref{eq:component-densities} builds on.

\paragraph{\gric.}
\gric addresses \emph{model-class} selection: it approximates the
evidence for manifolds of different dimensions and parameter counts
using a fixed-scale \mapsac/\msac data term plus complexity penalties
for the latent point coordinates and the global model parameters.  In
our setting the model class is fixed, so these penalties are constant
across hypotheses, while the data term still requires an externally
specified~$\sigma$.  Our contribution is complementary: starting from
the same codimension-based \gaussunif likelihood, we condition
on the candidate and its minimal sample, profile over the inlier
partition, and integrate~$\sigma^2$ exactly under an \invgamma
prior.  \gric approximates model-class evidence; our score removes the
inlier scale for within-class \ransac scoring.  Our $\nu = n_I\,d_g -
n_\theta$ counts residual Gaussian evidence after fitting
(\Cref{sec:supp-dof}); it is not \gric's complexity penalty.

\paragraph{The \magsacpp likelihood.}
\magsacpp~\cite{Barath2020} begins from a $\chi_4$ density because a
two-view correspondence has four coordinates.  The residual, however,
is a distance to the model manifold, so only the $d_g$ orthogonal
constraint directions remain ($d_g = 1$ for essential and fundamental
matrices, $d_g = 2$ for homographies): the relevant dimension is the
model codimension, not the ambient observation dimension
(\Cref{sec:supp-dims}).  More fundamentally, even $\chi_{d_g}$ is not
an appropriate model-ranking likelihood.  It is the density of the
\emph{radius} of a Gaussian residual vector and includes the volume
factor $r^{d_g-1}$; because $r = r(\btheta)$ depends on the candidate,
this is a parameter-dependent transformation of the observations, and
for $d_g \geq 2$ it assigns zero density to an exact fit.  Our
formulation instead starts from the Gaussian density of the
$d_g$-dimensional residual vector itself: $d_g$ is the per-observation
constraint dimension and $\nu$ the total residual evidence after
fitting, and neither is used as the shape of a radial likelihood.

\paragraph{Exact reparameterizations and a numerical equivalence.}
The threshold reparameterizations of \ransac, \msac and \gau
in~\cite{Shekhovtsov2025} are exact.  The reduction of \magsacpp to
fixed-scale \gau ($\sigma \approx 0.96\,\tau$ at observation dimension
$d = 4$)~\cite{Shekhovtsov2025} is a \emph{numerical} equivalence, not an
identity.  Our conclusions do not rely on that
approximation: our experiments run \magsacpp as defined in its paper
and reference code, and it nevertheless tracks \msac and \gau in
\Cref{fig:profile-grid,fig:sensitivity-grid}---in its optimum, its
degradation under perturbation, and its dependence on validation-set
size---which corroborates the equivalence empirically.
